\chapter*{Abstract}

Photovoltaic (PV) faults rarely produce visible symptoms and can only be inferred from deviations in measured electrical and meteorological signals, making automated Fault Detection and Diagnosis (FDD) essential as installations scale. This thesis develops an end-to-end AI-based FDD system, validated on real field data from the Costa PV plant and deployed on an embedded edge platform.

The system rests on a staged data pipeline (ingestion, temporal splitting, preprocessing, and feature engineering) that turns raw \mbox{1\,Hz} measurements into physics-informed features under strict leakage prevention: segment-aware temporal splitting and train-only statistic estimation at every stage. Two complementary tasks are then addressed. For anomaly detection, a semi-supervised detector is trained on normal operation alone, with a uniform Generalized Pareto Distribution policy calibrating thresholds identically across all detectors for fair comparison; the deployed physics-conditioned normalizing-flow model (PC-Flow, fewer than 25{,}000 parameters) attains a test binary PR-AUC of 0.989 and a worst-class PR-AUC of 0.968 under five-fold episode-stratified cross-validation. For fault classification, gradient-boosted trees with class weighting are used, the selected ExtraTrees classifier reaching a macro F1 of 0.959 and balanced accuracy of 94.3\% under five-fold cross-validation. The validated models are exported to ONNX and run on an NVIDIA Jetson Nano behind a real-time monitoring interface that turns the detector score into prioritised alarms, alongside per-class diagnosis and drift monitoring with operator-gated model updates.

Together these results trace a reproducible path from real-world PV measurements to reliable on-device fault awareness.

\vspace{1cm}
\noindent\textbf{Keywords:} photovoltaic systems, fault detection and diagnosis, machine learning, deep learning, predictive maintenance.
